\songcolumns{2}
	\beginsong{Prodavač - \trps}[by={Fešáci, Michal Tučný}]
	\newchords{Prodavac_verse\trps}
	\newchords{Prodavac_chorus\trps}
	\transpose{\trps}
	\beginverse*
		{\nolyrics \textbf{Intro:}} \[C]Pojďte \chordsoff všichni dovnitř, pozvěte si všechny známé,	
		my vám dobrou radu dáme, neboť právě otvíráme.
		Prodáváme, vyděláme, co kdo chcete, tak to máme,
		co nemáme, objednáme, všechno máme, všechno víme,
		poradíme, posloužíme.
		\chordson
	\endverse
	\beginverse\memorize[Prodavac_verse\trps]
		Stál \[C]krámek v naší ulici, v něm \[F]párky, buřty s hořčicí
		a \[G]bonbóny a sýr a sladký \[C]mák.\[G7]
		Tam \[C]chodíval jsem \[C7]potají, tak \[F]jak to kluci dělají
		a \[G]ochutnával od okurek \[C]lák.\[C7]
		A \[F]pro mou duši nevinnou pan \[C]vedoucí byl hrdinou,
		když \[D7]po obědě začal prodá\[G]vat.\[G7]
		Měl \[C]jazyk mrštný \[C7]jako bič a \[F]já byl z něho celý pryč
		a \[G]toužil jsem se prodavačem \[C]stát.
	\endverse
	\beginchorus\memorize[Prodavac_chorus\trps]
		\chordsoff
			\textbf{[Chorus]}
		\chordson
		\[C]Pět deka, deset deka, dvacet deka, třicet deka,
		\[F]kilo chleba, kilo cukru, jeden rohlík, jedna veka,
		\[G]všechno máme, co kdo chcete, obchod kvete, jen si račte \[C]říct.
		Čtyři kila, deset kilo, dvacet kilo, třicet kilo,
		\[F]navážíme, zabalíme, klaníme se, to by bylo,
		\[G]prosím pěkně, mohu nechat o jedenáct deka \[C]víc?
	\endchorus
	\beginverse\replay[Prodavac_verse\trps]
		Já ^nezapomněl na svůj cíl a ^záhy jsem se vyučil
		a ^moh být ze mě prodavačů ^král.^
		Je^nomže, jak ^běžel čas, náhle ^zaslechl jsem hudby hlas
		a z^nenadání na jevišti ^stál.^
		I když ^nejsem králem zpěváků, teď ^zpívám s partou Fešáků
		a ^nikdo vlastně neví, co jsem ^zač.^
		Mě ^potlesk hřeje ^do uší a ^mnohý divák netuší,
		^že mu vlastně zpívá proda^vač.
	\endverse
	\beginchorus\replay[Prodavac_chorus\trps]
		\chordsoff
			\textbf{[Chorus]}
		\chordson
		CHORUS
	\endchorus
	\beginverse\replay[Prodavac_verse\trps]
		Vím, ^že se život rozletí a ^sním o konci století,
		kdy ^nikdo neví, co je chvat a ^shon.^
    A dě^tem líčí ^babička, jak vy^padala elpíčka
    a ^co byl vlastně starý gramo^fón.^
    I kdy^by v roce dva tisíce, by^la veta po muzice,
    ob^chod je věc stále kvetou^cí.^
    Už ^se vidím, je ^to krása, ^ve výloze nápis hlásá
    ^Michal Tučný, odpovědný vedoucí.
	\endverse
	\beginchorus\replay[Prodavac_chorus\trps]
		\chordsoff
			\textbf{[Chorus]}
		\chordson
		CHORUS
	\endchorus
\endsong
